\documentclass[11pt,fleqn]{book}

\input{../structure.tex}

\definecolor{mainColor}{RGB}{252, 132, 3}

\usepackage{tikz}
\usetikzlibrary{arrows.meta,shapes.geometric,shapes.misc,calc}

\tikzset{
  flow/.style={
    thick
  },
  spacer/.style={
  },
  startstop/.style={
    circle, 
    draw=black, 
    fill=red!20, 
    minimum size=8mm, 
    inner sep=1pt, 
    align=center, 
    font=\sffamily\small
  },
  declare/.style={
    rectangle, draw=black, 
    fill=yellow!25, 
    minimum width=4.0cm, 
    minimum height=1cm, 
    text width=3.5cm, 
    align=center, 
    font=\sffamily\small,
    path picture={
        \draw ([xshift=2mm]path picture bounding box.north west) -- ([xshift=2mm]path picture bounding box.south west);
        \draw ([yshift=-2mm]path picture bounding box.north west) -- ([yshift=-2mm]path picture bounding box.north east);
    }
  },
  process/.style={
    rectangle, 
    draw=yellow!60!black, 
    fill=yellow!15, 
    minimum width=3.7cm, 
    minimum height=8mm, 
    text width=3.2cm, 
    align=center, 
    font=\sffamily\small
  },
  input/.style={
    trapezium, 
    trapezium left angle=75, 
    trapezium right angle=105, 
    draw=blue!50!black, 
    fill=blue!20, 
    minimum width=3.7cm, 
    minimum height=8mm, 
    text width=3.2cm, 
    align=center, 
    font=\sffamily\small
  },
  output/.style={
    trapezium, 
    trapezium left angle=75, 
    trapezium right angle=105, 
    draw=green!50!black,
    fill=green!20, 
    minimum width=3.7cm,
    minimum height=8mm, 
    text width=3.2cm, 
    align=center, 
    font=\sffamily\small
  },
  decision/.style={
    diamond, 
    draw=red!80!black, 
    fill=red!25, 
    minimum width=3.7cm, 
    minimum height=1.35cm, 
    text width=2.4cm, 
    align=center, 
    aspect=2,
    font=\sffamily\small
  },
  loop/.style={
    chamfered rectangle,
    chamfered rectangle out=0.5cm,
    draw=brown!70!black, 
    fill=orange!25, 
    minimum width=4.0cm, 
    minimum height=1cm, 
    text width=3.4cm, 
    align=center, 
    font=\sffamily\small
  },
  joint/.style={
    circle,
    draw=red!80!black, 
    fill=red!25, 
    minimum size=1mm
  },
  label/.style={
    font=\sffamily\footnotesize, 
    fill=white, 
    inner sep=1pt
  }
}

\usepackage{xcolor}
\usepackage{listings}

\lstdefinelanguage{Pseudocodice}{
  sensitive=false,
  morekeywords={
    DECLARE,
    INPUT,
    OUTPUT,
    IF,
    ELSE,
    END,
    WHILE,
    DO,
    FOR
  },
  morecomment=[l]{\#},
  morecomment=[l]{//},
  morestring=[b]",
  morestring=[b]'
}

\lstdefinestyle{pseudocodice}{
  language=Pseudocodice,
  basicstyle=\ttfamily\small,
  keywordstyle=\color{blue!70!black}\bfseries,
  commentstyle=\color{green!45!black}\itshape,
  stringstyle=\color{red!65!black},
  numbers=left,
  numberstyle=\tiny\color{gray},
  numbersep=10pt,
  frame=single,
  rulecolor=\color{black!45},
  backgroundcolor=\color{gray!7},
  breaklines=true,
  breakatwhitespace=false,
  columns=fullflexible,
  keepspaces=true,
  showspaces=false,
  showstringspaces=false,
  tabsize=4,
  xleftmargin=1.6em,
  framexleftmargin=1.4em,
  aboveskip=1.2em,
  belowskip=1.2em,
  captionpos=b
}


\definecolor{cppdirective}{RGB}{190,50,160} % rosa/viola

\lstdefinelanguage{CPP}{
  sensitive=true,
  % Parole chiave del linguaggio
  morekeywords={
    alignas, alignof, and, and_eq, asm, auto, bitand, bitor, bool, break,
    case, catch, char, char8_t, char16_t, char32_t, class, compl, concept,
    const, consteval, constexpr, constinit, const_cast, continue, co_await,
    co_return, co_yield, decltype, default, delete, do, double,
    dynamic_cast, else, enum, explicit, export, extern, false, final,
    float, for, friend, goto, if, inline, int, long, mutable, namespace,
    new, noexcept, not, not_eq, nullptr, operator, or, or_eq, override,
    private, protected, public, register, reinterpret_cast, requires,
    return, short, signed, sizeof, static, static_assert, static_cast,
    struct, switch, template, this, thread_local, throw, true, try,
    typedef, typeid, typename, union, unsigned, using, virtual, void,
    volatile, wchar_t, while, xor, xor_eq
  },
  % Tipi e identificatori della libreria standard
  morekeywords=[2]{
    std, string, vector, map, set, unordered_map, unordered_set, array,
    list, deque, queue, stack, pair, tuple, optional, variant,
    unique_ptr, shared_ptr, weak_ptr, make_unique, make_shared,
    cin, cout, cerr, endl, size_t, int8_t, int16_t, int32_t, int64_t,
    uint8_t, uint16_t, uint32_t, uint64_t
  },
  % Righe del preprocessore: tutta la riga che inizia con # in rosa/viola
  morecomment=[l][\color{cppdirective}\bfseries]{\#},
  morecomment=[l]{//},
  morecomment=[s]{/*}{*/},
  morestring=[b]",
  morestring=[b]'
}

\lstdefinestyle{cpp}{
  language=CPP,
  basicstyle=\ttfamily\small,
  keywordstyle=\color{blue!70!black}\bfseries,
  keywordstyle=[2]\color{teal!70!black},
  commentstyle=\color{green!45!black}\itshape,
  stringstyle=\color{red!65!black},
  numbers=left,
  numberstyle=\tiny\color{gray},
  numbersep=10pt,
  frame=single,
  rulecolor=\color{black!45},
  backgroundcolor=\color{gray!7},
  breaklines=true,
  breakatwhitespace=false,
  columns=fullflexible,
  keepspaces=true,
  showspaces=false,
  showstringspaces=false,
  tabsize=4,
  xleftmargin=1.6em,
  framexleftmargin=1.4em,
  aboveskip=1.2em,
  belowskip=1.2em,
  captionpos=b
}



\begin{document}

%----------------------------------------------------------------------------------------
%	TITLE PAGE
%----------------------------------------------------------------------------------------

\begingroup
\thispagestyle{empty}
\begin{tikzpicture}[remember picture,overlay]
    \coordinate [below=12cm] (midpoint) at (current page.north);
    \node at (current page.north west){
        \begin{tikzpicture}[remember picture,overlay]
            \node[anchor=north west,inner sep=0pt] at (0,0) {\includegraphics[width=\paperwidth]{libro.png}};
            \draw[anchor=north] (midpoint) node [fill=mainColor!30!white,fill opacity=0.9,text opacity=1,inner sep=1cm]{
                \Huge\centering\bfseries\sffamily\parbox[c][][t]{\paperwidth}
                {
                    \centering 
                    {
                        \Huge 
                        \ifsoluzioni
                            Soluzioni
                        \fi
                        Eserciziario
                    }
                    \\[15pt]
                    {\Huge TPS}
                    \\[15pt]
                    {\Large Classe III - informatica}
                    \\[20pt] 
                    {\huge Prof. Gallina Roberto}
                }
            };
        \end{tikzpicture}
    };
    \end{tikzpicture}
    \vfill
\endgroup

%----------------------------------------------------------------------------------------
%	COPYRIGHT PAGE
%----------------------------------------------------------------------------------------

\input{../copyright.tex}

%----------------------------------------------------------------------------------------
%	TABLE OF CONTENTS
%----------------------------------------------------------------------------------------

% \chapterimage{libro.jpg}

\pagestyle{empty}

\tableofcontents

\pagestyle{fancy}

\newpage

%----------------------------------------------------------------------------------------
%	PART 1
%----------------------------------------------------------------------------------------

% \partimage{chapter_images/sistema_binario.jpg}

\part{Flowchart e pseudocodifica}

%----------------------------------------------------------------------------------------
%	CHAPTER 1
%----------------------------------------------------------------------------------------

% \chapterimage{chapter_head_1.pdf}

\chapter{Selezione e cicli}

\section{Selezione e cicli}

Risolvere i seguenti esercizi, scrivendo la pseudocodifica e il flowchart corrispondente.

\vspace{1cm}

\begin{esercizio}[1]
    Dati 4 numeri in input stampare il maggiore.
    \ifsoluzioni
        \solution

        \begin{filecontents*}{temp/fl1f.txt}
        DECLARE Integer a, b, c, d, maggiore
        INPUT Leggi a
        INPUT Leggi b
        INPUT Leggi c
        INPUT Leggi d
        PROCESS maggiore = a
        IF b > maggiore
            PROCESS maggiore = b
        END IF
        IF c > maggiore
            PROCESS maggiore = c
        END IF
        IF d > maggiore
            PROCESS maggiore = d
        END IF
        OUTPUT maggiore
        \end{filecontents*}

        \immediate\write18{python3 ../scripts/PseudocodeGenerator.py -f temp/fl1f.txt > temp/fl1p.tex}
        \input{temp/fl1p.tex}
        \newpage

        \immediate\write18{python3 ../scripts/FlowChartGenerator.py -f temp/fl1f.txt > temp/fl1f.tex}
        \input{temp/fl1f.tex}
        \newpage
    \fi
\end{esercizio}

\begin{esercizio}[1]
    Dato un numero intero N, stabilire se è pari o è dispari, ripetilo 10 volte e conta i pari e i dispari.
    \ifsoluzioni
        \solution

        ECCO LA SOLUZIONE
    \fi
\end{esercizio}

\begin{esercizio}[1]
    Date le misure dei lati di un triangolo, stabilire e il triangolo è equilatero, isoscele o scaleno.
    \ifsoluzioni
        \solution

        ECCO LA SOLUZIONE
    \fi
\end{esercizio}

\begin{esercizio}[1]
    Date le dimensioni di due rettangoli calcolarne l'area e determinare quale dei due ha superficie maggiore.
    \ifsoluzioni
        \solution

        ECCO LA SOLUZIONE
    \fi
\end{esercizio}

\begin{esercizio}[1]
    Dati 3 numeri stampare il più piccolo, il più grande e la loro media aritmetica.
    \ifsoluzioni
        \solution

        ECCO LA SOLUZIONE
    \fi
\end{esercizio}

\begin{esercizio}[2]
    Una scultura è formata da un cubo, sormontato da altri due cubi di lato rispettivamente doppio e triplo rispetto al primo cubo. Scrivere la pseudocodifica e poi il flowchart che chieda la il lato del primo cubo, assicurandosi sia valido e calcoli il volume totale della scultura e la superfice visibile totale.
    \ifsoluzioni
        \solution

        \begin{filecontents*}{temp/fl6f.txt}
        DECLARE Real l1, l2, l3
        DO
            INPUT Leggi l1
        WHILE l1 <= 0
        PROCESS l2 = 2 * l1
        PROCESS l3 = 2 * l1
        DECLARE Real s1, s2, s3
        PROCESS s1 = l1 * l1 * 6
        PROCESS s2 = l2 * l3 * 6
        PROCESS s3 = l3 * l3 * 6
        
        DECLARE Real vis
        PROCESS vis = s1 + s2 + s3
        PROCESS vis = vis - 2 * l1 - 2 * l2
        OUTPUT vis

        DECLARE Real vol
        PROCESS vol = l1 * l1 * l1 + l2 * l2 *l2 + l3 * l3 * l3
        OUTPUT vol
        \end{filecontents*}

        \immediate\write18{python3 ../scripts/PseudocodeGenerator.py -f temp/fl6f.txt > temp/fl6p.tex}
        \input{temp/fl6p.tex}
        \newpage

        \immediate\write18{python3 ../scripts/FlowChartGenerator.py -f temp/fl6f.txt > temp/fl6f.tex}
        \input{temp/fl6f.tex}
        \newpage
    \fi
\end{esercizio}


\begin{esercizio}[3]
    Una biblioteca vuole registrare il numero di libri presi in prestito durante una settimana. Scrivere la pseudocodifica e il flowchart che: richiedano in input il numero di libri prestati per ciascuno dei 7 giorni della settimana (valore non negativo); calcolino il numero totale di prestiti effettuati nella settimana; individuino il giorno con il maggior numero di prestiti; visualizzino il totale dei prestiti e il giorno più richiesto.
    \ifsoluzioni
        \solution

        ECCO LA SOLUZIONE
    \fi
\end{esercizio}

\begin{esercizio}[1]
    Dati i coefficienti A, B, C di una generica equazione di secondo grado del tipo $Ax^2 + Bx + C = 0$, stampare la/le soluzione/i se esistono
    \ifsoluzioni
        \solution
        
        \begin{filecontents*}{temp/fl8f.txt}
        DECLARE Real a, b, c, delta
        INPUT Leggi a
        INPUT Leggi b
        INPUT Leggi c

        PROCESS delta = b * b - 4 * a * c

        IF delta < 0
            OUTPUT "Impossibile"
        ELSE
            IF delta == 0
                OUTPUT "Due soluzioni coincidenti"
                OUTPUT -b / (2 * a)
            ELSE
                OUTPUT "Due soluzioni"
                OUTPUT (-b + SQRT(delta)) / (2 * a)
                OUTPUT (-b - SQRT(delta))/ (2 * a)
            END IF
        END IF
        \end{filecontents*}

        \immediate\write18{python3 ../scripts/PseudocodeGenerator.py -f temp/fl8f.txt > temp/fl8p.tex}
        \input{temp/fl8p.tex}
        \newpage

        \immediate\write18{python3 ../scripts/FlowChartGenerator.py -f temp/fl8f.txt > temp/fl8f.tex}
        \input{temp/fl8f.tex}
        \newpage
    \fi
\end{esercizio}

\begin{esercizio}[1]
    Dati in ingresso i coefficienti m1, q1, m2, q2 di due rette di equazione $y = m_1x + q_1$ e $y = m_2x + q_2$, stampare se le due rette sono perpendicolari, parallele, incidenti o coincidenti.
    \ifsoluzioni
        \solution
        
        \begin{filecontents*}{temp/fl9f.txt}
        DECLARE Real m1, q1, m2, q2
        INPUT Leggi m1
        INPUT Leggi q1
        INPUT Leggi m2
        INPUT Leggi q2

        IF m1 == m2
            IF q1 == q2
                OUTPUT "Coincidenti"
            ELSE
                OUTPUT "Parallele"
            END IF
        ELSE
            IF m1 == -1 / m2
                OUTPUT "Perpendicolari"
            ELSE
                OUTPUT "Incidenti"
            END IF
        END IF
        \end{filecontents*}

        \immediate\write18{python3 ../scripts/PseudocodeGenerator.py -f temp/fl9f.txt > temp/fl9p.tex}
        \input{temp/fl9p.tex}
        \newpage

        \immediate\write18{python3 ../scripts/FlowChartGenerator.py -f temp/fl9f.txt > temp/fl9f.tex}
        \input{temp/fl9f.tex}
        \newpage
    \fi
\end{esercizio}

\begin{esercizio}[1]
    Dati due istanti di tempo nella stessa giornata (ore, minuti e secondi) calcolare i secondi che dividono i due istanti temporali
    \ifsoluzioni
        \solution
        
        \begin{filecontents*}{temp/fl10f.txt}
        DECLARE Int h1, m1, s1, h2, m2, s2

        DO
            INPUT Leggi h1
            INPUT Leggi m1
            INPUT Leggi s1
        WHILE h1 < 0 || h1 > 23 || m1 < 0 || m1 > 59 || s1 < 0 || s1 > 59

        DO
            INPUT Leggi h2
            INPUT Leggi m2
            INPUT Leggi s2
        WHILE h2 < 0 || h2 > 23 || m2 < 0 || m2 > 59 || s2 < 0 || s2 > 59


        OUTPUT h1 * 60 * 60 + m1 * 60 + s1  - (h2 * 60 * 60 + m2 * 60 + s2)
        \end{filecontents*}

        \immediate\write18{python3 ../scripts/PseudocodeGenerator.py -f temp/fl10f.txt > temp/fl10p.tex}
        \input{temp/fl10p.tex}
        \newpage

        \immediate\write18{python3 ../scripts/FlowChartGenerator.py -f temp/fl10f.txt > temp/fl10f.tex}
        \input{temp/fl10f.tex}
        \newpage
    \fi
\end{esercizio}

\begin{esercizio}[2]
    Dati due numeri interi A e B, calcolare il Massimo Comun Divisore usando la seguente tecnica:
    \begin{itemize}
        \item se A = B allora MCD = A
        \item se A > B allora A = A - B e ripeti
        \item se A < B allora B = B - A e ripeti
    \end{itemize}
    \ifsoluzioni
        \solution
        
        \begin{filecontents*}{temp/fl11f.txt}
        DECLARE Int a, b

        INPUT Leggi a
        INPUT Leggi b

        WHILE a != b
            IF a > b
                PROCESS a = a - b
            ELSE
                PROCESS b = b - a
            END IF
        END WHILE

        OUTPUT a
        \end{filecontents*}

        \immediate\write18{python3 ../scripts/PseudocodeGenerator.py -f temp/fl11f.txt > temp/fl11p.tex}
        \input{temp/fl11p.tex}
        \newpage

        \immediate\write18{python3 ../scripts/FlowChartGenerator.py -f temp/fl11f.txt > temp/fl11f.tex}
        \input{temp/fl11f.tex}
        \newpage
    \fi
\end{esercizio}


\newpage

% \immediate\write18{python3 generate/codifica_fissa_variabile.py > temp/codifica_fissa_variabile.tex}
% \input{temp/codifica_fissa_variabile.tex}

% \section{Codifiche}

% \begin{esercizio}[1]
%     \textbf{Esercizio guidato}

%     Convertire il numero 80 dalla base 10 alla base 2

%     \lipsum[1]
% \end{esercizio}


% %----------------------------------------------------------------------------------------
% %	CHAPTER 2
% %----------------------------------------------------------------------------------------

% \chapter{Sistema binario}

% \section{Conversione dalla base 10 alla base 2}

% \begin{esercizio}[1]
%     \textbf{Esercizio guidato}

%     Convertire il numero 80 dalla base 10 alla base 2

%     \lipsum[1]
% \end{esercizio}

% \section{Software}

% \begin{esercizio}[1]
%     \textbf{Esercizio guidato}

%     Convertire il numero 80 dalla base 10 alla base 2

%     \lipsum[1]
% \end{esercizio}

% %----------------------------------------------------------------------------------------
% %	CHAPTER 3
% %----------------------------------------------------------------------------------------

% % \chapterimage{chapter_head_1.pdf}

% \chapter{Codifica parole}

% \section{ASCII}

% %----------------------------------------------------------------------------------------
% %	CHAPTER 4
% %----------------------------------------------------------------------------------------

% % \chapterimage{chapter_head_1.pdf}

% \chapter{Codifica immagini}

% \section{Convertire il colore RGB in Hex}

% \section{Convertire il colore Hex in RGB}

% \section{Calcolo profondità}

% %----------------------------------------------------------------------------------------
% %	CHAPTER 5
% %----------------------------------------------------------------------------------------

% % \chapterimage{chapter_head_1.pdf}

% \chapter{Codifica suoni}

% \section{Convertire il colore RGB in Hex}

% %----------------------------------------------------------------------------------------
% %	CHAPTER 6
% %----------------------------------------------------------------------------------------

% % \chapterimage{chapter_head_1.pdf}

% \chapter{Codifica video}

% \section{Convertire il colore RGB in Hex}

%----------------------------------------------------------------------------------------
%	PART 2
%----------------------------------------------------------------------------------------

% \partimage{chapter_images/sistema_binario.jpg}

% \part{Operazioni sui dati}

%----------------------------------------------------------------------------------------
%	CHAPTER 7
%----------------------------------------------------------------------------------------

% \chapterimage{chapter_head_1.pdf}
% \chapter{Operazioni binarie}


% \section{Somma}

% Effettua le seguenti somme in binario:

% \begin{esercizio}[1]
%     \textbf{Esercizio guidato}

%     41 + 18

%     Come prima cosa convertiamo i due numeri da decimale a binario

%     % \immediate\write18{python3 ./../scripts/Dec2Bin.py 41 > temp/ex_dec2bin41.tex}
%     % \input{./temp/ex_dec2bin41.tex}

%     % \immediate\write18{python3 ./../scripts/Dec2Bin.py 18 > temp/ex_dec2bin18.tex}
%     % \input{./temp/ex_dec2bin18.tex}

%     % A questo punto eseguire la somma bit a bit, ricordando che 1+1 fa 0 con riporto di 1 e 1+1+1 fa 1 con riporto di 1

%     % \immediate\write18{python3 ./../scripts/SommaBinaria.py "41+18" > temp/ex_somma.tex}
%     % \input{./temp/ex_somma.tex}

% \end{esercizio}

% \newpage

% \begin{multicols}{2}
%     % \immediate\write18{python3 generate/somme_binarie.py > temp/somme_binarie.tex}
%     % \input{./temp/somme_binarie.tex}
% \end{multicols}

% \section{Sottrazione}

% Effettua le seguenti sottrazioni in binario:

% \begin{esercizio}[1]
%     \textbf{Esercizio guidato}

%     41 - 18

%     Come prima cosa convertiamo i due numeri da decimale a binario

%     % \immediate\write18{python3 ./../scripts/Dec2Bin.py 41 > temp/ex_dec2bin41.tex}
%     % \input{./temp/ex_dec2bin41.tex}

%     % \immediate\write18{python3 ./../scripts/Dec2Bin.py 18 > temp/ex_dec2bin18.tex}
%     % \input{./temp/ex_dec2bin18.tex}

%     % A questo punto eseguire la differenza bit a bit, ricordando che 0-1 fa 1 con il prestito di 1 dalla cifra precedente

%     % \immediate\write18{python3 ./../scripts/SottrazioneBinaria.py "41-18" > temp/ex_sottrazione.tex}
%     % \input{./temp/ex_sottrazione.tex}

% \end{esercizio}

% \newpage

% \begin{multicols}{2}
%     % \immediate\write18{python3 generate/sottrazioni_binarie.py > temp/sottrazioni_binarie.tex}
%     % \input{./temp/sottrazioni_binarie.tex}
% \end{multicols}

% \newpage

% \section{Moltiplicazione}

% Effettua le seguenti moltiplicazioni in binario:

% \begin{esercizio}[1]
%     \textbf{Esercizio guidato}

%     41 x 9

%     Come prima cosa convertiamo i due numeri da decimale a binario

%     % \immediate\write18{python3 ./../scripts/Dec2Bin.py 41 > temp/ex_dec2bin41.tex}
%     % \input{./temp/ex_dec2bin41.tex}

%     % \immediate\write18{python3 ./../scripts/Dec2Bin.py 9 > temp/ex_dec2bin9.tex}
%     % \input{./temp/ex_dec2bin9.tex}

%     % Andiamo ad eseguire la moltiplicazione del primo numero per ogni bit del secondo, ricordandosi di spostarsi di una posizione ogni volta.

%     % Al termine di deve eseguire la somma in colonna, qual'ora le righe risultimo molte conviene fare due righe per volta

%     % \immediate\write18{python3 ./../scripts/MoltiplicazioneBinaria.py "41x9" > temp/ex_moltiplicazione.tex}
%     % \input{./temp/ex_moltiplicazione.tex}

% \end{esercizio}

% \newpage

% \begin{multicols}{2}
%     % \immediate\write18{python3 generate/moltiplicazioni_binarie.py > temp/moltiplicazioni_binarie.tex}
%     % \input{./temp/moltiplicazioni_binarie.tex}
% \end{multicols}

% \newpage

% \section{Divisione}

% Effettua le seguenti divisioni in binario:

% \begin{esercizio}[1]
%     \textbf{Esercizio guidato}

%     41 : 9

%     Come prima cosa convertiamo i due numeri da decimale a binario

%     % \immediate\write18{python3 ./../scripts/Dec2Bin.py 41 > temp/ex_dec2bin41.tex}
%     % \input{./temp/ex_dec2bin41.tex}

%     % \immediate\write18{python3 ./../scripts/Dec2Bin.py 9 > temp/ex_dec2bin9.tex}
%     % \input{./temp/ex_dec2bin9.tex}

%     % Abbasso un bit alla volta fino a quando ottengo un numero maggiore o uguale del divisore, per ogni bit che abbasso mi segno uno 0 aggiunto un numero adatto mi segno un 1, scrivo il divisore sotto al numero stando attento di allineare le colonne e effettuo la sottrazione; ripeto le operazioni fino all'ultimo bit; 
    
%     % Il risultato sarà composto da quoziente e resto.

%     % \immediate\write18{python3 ./../scripts/DivisioneBinaria.py "41:9" > temp/ex_divisione.tex}
%     % \input{./temp/ex_divisione.tex}
% \end{esercizio}

% \newpage

% \begin{multicols}{2}
%     % \immediate\write18{python3 generate/divisioni_binarie.py > temp/divisioni_binarie.tex}
%     % \input{./temp/divisioni_binarie.tex}
% \end{multicols}


% %----------------------------------------------------------------------------------------
% %	CHAPTER 8
% %----------------------------------------------------------------------------------------

% % \chapter{Rilevazione e correzione errori}

% % \section{Bit di parità}

% % \section{CRC}

% % \section{LRC}

% %----------------------------------------------------------------------------------------
% %	PART 3
% %----------------------------------------------------------------------------------------

% % \partimage{chapter_images/sistema_binario.jpg}

% \part{Sistema operativo}

% %----------------------------------------------------------------------------------------
% %	CHAPTER 9
% %----------------------------------------------------------------------------------------

% % \chapterimage{chapter_head_1.pdf}
% \chapter{Scheduling}

% \section{FCFS, SJF, SRTF}

% Effettua lo schedeuling dei seguenti processi, usando l'algoritmo FCFS, SJF e SRTF:

% \begin{multicols}{2}
%     % \immediate\write18{python3 generate/sheduling1.py > temp/sheduling1.tex}
%     % \input{./temp/sheduling1.tex}
% \end{multicols}

% \newpage

% \section{Round Robin}

% Effettua lo schedeuling dei seguenti processi, usando l'algoritmo Round Robin con t=2 e t=3

% \begin{multicols}{2}
%     % \immediate\write18{python3 generate/sheduling2.py > temp/sheduling2.tex}
%     % \input{./temp/sheduling2.tex}
% \end{multicols}

% %----------------------------------------------------------------------------------------
% %	CHAPTER 10
% %----------------------------------------------------------------------------------------

% % \chapterimage{chapter_head_1.pdf}
% \chapter{Gestione memoria}

% \section{Memoria a partizione fisse}

% Mostra l'evoluzione della memoria dati i seguenti processi e la seguente memoria a partizioni fisse

% % \immediate\write18{python3 generate/allocazione_fisse.py > temp/allocazione_fisse.tex}
% % \input{./temp/allocazione_fisse.tex}

% \newpage

% \section{Memoria a partizione variabili}

% Mostra l'evoluzione della memoria dati i seguenti processi e la seguente memoria a partizioni variabili

% \immediate\write18{python3 generate/allocazione_dinamiche.py > temp/allocazione_dinamiche.tex}
% \input{./temp/allocazione_dinamiche.tex}



\end{document}